%% paginas/1_externa/lombada.tex
%% Elemento Externo Opcional (Lombada) -- Conforme ABNT NBR 12225
%% NOTA: Útil para trabalhos impressos encadernados (capa dura).
%% Exige o pacote 'graphicx' se for usar rotação de texto.

\begin{titlepage}
  \thispagestyle{empty}
  \begin{center}
    
    \vspace*{2cm}
    
    \begin{spacing}{1.0}
      \noindent
      \textbf{LOMBADA} \\
      \vspace{1cm}
      Este arquivo é um modelo de Lombada acadêmica. \\
      Geralmente contém as informações dispostas longitudinalmente para impressão na lombada do livro:
      
      \vspace{2cm}
      
      % Exemplo de formatação longitudinal:
      {\large\MakeUppercase{\meuautor}} \hfill {\large\MakeUppercase{\meutitulo}} \hfill {\large\meuano}
      
      \vspace{4cm}
      
      \begin{adjustwidth}{2cm}{2cm}
        \small
        \textit{Dica: Se for imprimir em capa dura, a gráfica normalmente gera a lombada com base na espessura do volume de folhas (número de páginas). Este modelo serve como um guia conceitual.}
      \end{adjustwidth}
      
    \end{spacing}
    
  \end{center}
\end{titlepage}
